\paragraph{Conditional Moran quartet image.}
Consider the finite-$M$ Moran-integrated MSRC recursion with a fixed rooted
four-taxon species tree and condition on the observed arrangement pattern
$Z=z$. For the neutral, zero-switching subfamily with equal branch copy
counts, point-mass initial count $K\in\{1,\ldots,M-1\}$, and zero lower
branch durations, let $T$ be the root ancestral duration. For
$0011$, $0101$, and $0110$, respectively, the compatible pair determines
$Q_1=12|34$, $Q_2=13|24$, and $Q_3=14|23$ when it coalesces before the root.

\begin{proposition}[Vertex closure]
For either $((1,2),(3,4))$ or $(((1,2),3),4)$, the conditional quartet law
for the corresponding pattern satisfies
\[
q^{(j)}(T)=(1-S_j(T))e_j+S_j(T)r_j(T),
\]
where $S_j(T)$ is the compatible-pair survival probability and $r_j(T)$ is a
probability vector. Moreover $S_j(T)\to0$ as $T\to\infty$, the conditioning
event has positive probability for every finite $T$, and hence
\[
e_j\in\overline{\{q(z_j;\theta):P_\theta(Z=z_j)>0\}}.
\]
\end{proposition}

\begin{proof}
With zero switching and zero lower-branch durations, the compatible pair in
pattern $z_j$ enters the first ancestral branch in one background. At every
interior Moran count its pair-coalescence rate is positive. Thus its survival
probability to the root is a finite-branch survival probability $S_j(T)$
that tends to zero with the root duration. On the complementary event the
first compatible coalescence fixes the unrooted quartet to $Q_j$, regardless
of the later placement of the fourth lineage. The unresolved event contributes
a probability vector $r_j(T)$. The joint Moran recursion retains the sampling
mass of $z_j$; since $0<K<M$, that mass is positive, so division by it is
valid. Therefore $q^{(j)}(T)\to e_j$. The pectinate case has the same
argument because the sister and intermediate branches have zero duration.
\end{proof}

This proposition proves vertex closure only. It does not prove that the
conditional image contains a sector, that it is convex, or that its union
over patterns is dense in the full simplex. Establishing any of those claims
requires a separate global inverse, degree, mixture, or obstruction theorem.
